\section{Related work}
\label{sec:related}

\paragraph{GPU state machines and device-side control.}
FLAME GPU 2 maps agent memory, state transitions, population partitioning, and
communication to GPUs \cite{richmond2023flamegpu2}. CUDA Graphs provide replay,
conditional nodes, and device graph launch \cite{nvidia2026cudaprogrammingguide}.
GPUOS uses a persistent GPU worker kernel and atomic queues fed from a
host-managed work queue \cite{yang2026gpuos}; MPK and
Event Tensor compile tensor task graphs into persistent megakernels
\cite{cheng2026mpk,jin2026eventtensor}. These systems provide mechanisms for
device-side control and launch amortization. Our measurements ask a different
question: whether agent traces supply enough work under a declared grouping
before launch deadlines, and what one matched host observation costs.

\paragraph{Agent and compound-AI serving.}
A production and open-source characterization by Yang et al. identifies
repeated CPU--GPU crossings, host orchestration on the critical path, and
bursty CPU demand in agentic workflows; its Agora prototype pools and schedules
CPU resources while oversubscribing GPU memory
\cite{yang2026agenticarchitecture}. ThunderAgent represents workflows as LLM
Programs and jointly manages KV caches, system state, and tool environments
with a program-aware scheduler \cite{kang2026thunderagent}.
Parrot exposes semantic variables and application dataflow for cross-request
optimization \cite{lin2024parrot}. Agentix schedules agent programs using
execution progress \cite{luo2026agentix}; SAGA uses KV-cache-aware execution
graphs and session-affinity batching \cite{guo2026saga}; and InferCept manages
inference state across pauses for external interactions
\cite{abhyankar2024infercept}. Murakkab jointly maps declarative workflow
components to models and hardware under SLOs \cite{chaudhry2026murakkab}. The
MARS preprint coordinates GPU inference and CPU tool execution through an
external control plane \cite{wang2026mars}; the concurrent OpRAG preprint uses
persistent workers, bounded queues, and batched GPU operators for multi-stage
RAG \cite{sarker2026oprag}.
AgentServe co-schedules multi-agent inference on one GPU
\cite{zhang2026agentserve}. These systems operate at model-call, tool, or
workflow granularity. Our unit is the deterministic transition after a model or
tool event has completed.

\paragraph{Non-model CPU and GPU placement.}
Agentic CPU-GPU Scheduling profiles a complete AI tool and selects immediate
GPU, queued GPU, or CPU execution \cite{lu2026agenticcpugpu}. TokTier moves
stateful tokenization across CPU and GPU while enforcing exact token identity
\cite{zhang2026toktier}. Measurements of CPU-induced LLM slowdowns show that
CPU starvation can delay launches, communication, and tokenization even with
CUDA Graphs, and that adding CPU cores can be the economical remedy
\cite{chung2026cpuslowdowns}. These are the closest tool-level and methodological
comparisons. They make a tuned CPU implementation and CPU provisioning
mandatory baselines for the proposed online runtime.

\paragraph{Batching and clustering.}
Classical real-time scheduling studies job families, time windows, and
compatible batching \cite{barnoy2009throughput}. Joint batching work addresses
heterogeneous asynchronous arrivals and deadlines \cite{cang2024jointbatching},
while serial-batch scheduling includes minimum batch size, job families, and
release times \cite{huertas2025serialbatch}. One-dimensional $r$-gathering
studies minimum-occupancy clusters, outliers, contiguous solutions, and exact
line algorithms \cite{aggarwal2010anonymity,akagi2015rgather,
nakano2019rgather,sarker2021rgather}. Univariate microaggregation with
suppression supplies another dynamic-programming analogue
\cite{laszlo2013microaggregation}. GPU dynamic batching has also been modeled
with batch-dependent service, stochastic arrivals, latency, and power
\cite{xu2023smdpbatching}. The recurrence in this paper is a restricted trace
oracle inside that established territory, not an algorithmic-priority claim.

\paragraph{Positioning.}
\claimid{RC-012}
Prior agent systems characterize serving, scheduling, state management, and
tool execution, while batching theory characterizes which arrivals can be
served together under timing and occupancy constraints. The supply of
route-key-conditioned post-event transitions above a hardware crossover lies
at their interface. Observation placement adds the second axis: a cohort may
exist, yet a host round trip can still dominate a short control chain. This
paper gives those axes common boundary quantities and measures each with a
separate, auditable experiment.
